
\begin{frame}
\frametitle{Cartucho de Recuperación (REC) (1)}
\begin{itemize}
\item Estudiante tiene derecho a solicitar SOLAMENTE UNA VEZ un proyecto de recuperación aplicable a un solo proyecto o actividad durante todo el cuatrimestre. 
\item Esta solicitud debe HACERLA es estudiante - El profesor NO ES RESPONSABLE de informar al estudiante cuando tiene un ADEUDO. Cuando el estudiante realice dicho acercamiento, es responsabilidad que haya un respaldo del proyecto de recuperación, incluyendo fechas de entrega.
\item A partir de la primera solicitud de recuperación de un estudiante, se asigna una fecha de entrega. Este tiempo corre para todos aquellos estudiantes que tengan adeudo, por lo que si otros estudiantes se acercan de manerá tardía, tendrán el mismo tiempo para la entrega del proyecto.
\item Si alguien hace la solicitud de recuperación posterior a la fecha de entrega asignada inicialmente, el proyecto \textbf{YA NO SE PUEDE RECUPERAR.}
\end{itemize}
\end{frame}

\begin{frame}
\frametitle{Cartucho de Recuperación (REC) (2)}
\begin{itemize}
\item Si el proyecto no entregado es individual, se asigna otro proyecto diferente.
\item Si el proyecto es en equipo, de común acuerdo con los integrantes pueden trabajar en otro proyecto diferente en equipo, o recibir una asignación individual de un proyecto diferente.
\item El proyecto de recuperación será asignado siempre y cuando el estudiante cumpla lo siguiente:
\begin{itemize}
\item El estudiante debe tener el porcentaje de asistencia mínimo necesario para aprobar. 
\item El estudiante ha agendado el porcentaje de asesorias. 
\end{itemize}
\item El nuevo proyecto asignado esta diseñado para que el estudiante invierta en él por lo menos 2 SEMANA. Si lo solicita un día antes de terminar el cuatrimestre, posiblemente no tendrá tiempo de llevarlo a cabo. 

\end{itemize}
\end{frame}

